% --- Archon physics formalization source begin ---
% archon:physics
% archon:covers PhyXMiniProblems/problem_phyx_mini_0586.lean
% archon:source-report reports/phyx_mini/problem_phyx_mini_0586.source.json

\chapter{Physics problem phyx\_mini\_0586}
\label{ch:PhyXMiniProblems_problem_phyx_mini_0586}

\paragraph{Problem source.}
\#\# Physical scenario

Figure shows the energy-level diagram for a finite, one-dimensional energy well that contains an electron. The nonquantized region begins at \textbackslash{}( E\_4 = 450.0\textbackslash{},\textbackslash{}text\{eV\} \textbackslash{}).The electron can absorb at the indicated wavelengths: \textbackslash{}( \textbackslash{}lambda\_a = 14.588\textbackslash{},\textbackslash{}text\{nm\} \textbackslash{}) and \textbackslash{}( \textbackslash{}lambda\_b = 4.8437\textbackslash{},\textbackslash{}text\{nm\} \textbackslash{}) and for any wavelength less than \textbackslash{}( \textbackslash{}lambda\_c = 2.9108\textbackslash{},\textbackslash{}text\{nm\} \textbackslash{}).

\#\# Question

What is the energy of the first excited state?

\#\# Answer choices

- A: A: 85eV

- B: B: 24eV

- C: C: 109eV

- D: D: 426eV

\#\# Figure caption (auxiliary; use the image as primary evidence)

The image is an energy level diagram. It shows quantized energy levels represented by horizontal lines, labeled \textbackslash{}(E\_1\textbackslash{}), \textbackslash{}(E\_2\textbackslash{}), \textbackslash{}(E\_3\textbackslash{}), and \textbackslash{}(E\_4\textbackslash{}). These levels are evenly spaced along the vertical axis labeled "Energy." Above the highest quantized level \textbackslash{}(E\_4\textbackslash{}), there is a shaded region labeled "Nonquantized." The energy axis starts from zero at the bottom. The arrangement suggests discrete energy levels up to \textbackslash{}(E\_4\textbackslash{}) and a continuum above it.

\#\# Dataset metadata

Quantum Phenomena; Physical Model Grounding Reasoning, Multi-Formula Reasoning

\paragraph{Recorded answer/context.}
C: C: 109eV

\paragraph{Figure/image path.}
/root/proposal\_for\_physic/science-mango/phyx\_mini\_run/phyx\_data/test\_image/586.png

\paragraph{Formalization target.}
create a compiling Lean file with sorry bodies at `PhyXMiniProblems/problem\_phyx\_mini\_0586.lean`. The Lean declarations must preserve the physical quantities, dimensions or dimensional roles, figure labels, governing-law hypotheses, and final relation expressed by this problem.
Use Mathlib/Physlib names found through LeanExplore where available. If a physics API is missing, introduce faithful local abstractions rather than scalar placeholder aliases.

\begin{theorem}[Physics formalization target]
\label{thm:physics:phyx_mini_0586:target}
\lean{PhyXMiniProblems.ProblemPhyXMini0586.problem_phyx_mini_0586}
\uses{def:physics:phyx-mini-0586:phyxminiproblems-problemphyxmini0586-energyinelectronvolts, def:physics:phyx-mini-0586:phyxminiproblems-problemphyxmini0586-photonenergywavelengthconstantevnm, def:physics:phyx-mini-0586:phyxminiproblems-problemphyxmini0586-finitewellabsorptionsetup, def:physics:phyx-mini-0586:phyxminiproblems-problemphyxmini0586-matchesfiniteonedimensionalwellscenario, def:physics:phyx-mini-0586:phyxminiproblems-problemphyxmini0586-matchesproblemreadouts, def:physics:phyx-mini-0586:phyxminiproblems-problemphyxmini0586-matchesprimaryenergylevelfigure, def:physics:phyx-mini-0586:phyxminiproblems-problemphyxmini0586-hasphysicalfinitewellparameters, def:physics:phyx-mini-0586:phyxminiproblems-problemphyxmini0586-satisfiesfinitewellphotonabsorptionlaws, def:physics:phyx-mini-0586:phyxminiproblems-problemphyxmini0586-recordeddatasetanswer, def:physics:phyx-mini-0586:phyxminiproblems-problemphyxmini0586-roundstonearestelectronvolt, def:physics:phyx-mini-0586:phyxminiproblems-problemphyxmini0586-isuniquematchingfirstexcitedenergy}
The assigned autoformalize agent should translate this physics problem into Lean declarations in the covered file, with theorem and lemma proof bodies written as `by sorry`.
\end{theorem}
\begin{proof}
This is an autoformalization task, not a proof task. Produce statements that can later be proved by the physics prover without weakening the physical model.
\end{proof}

% --- Archon declaration topology begin ---
\section{Declaration topology}
\begin{definition}[Energy Quantity]
  \label{def:physics:phyx-mini-0586:phyxminiproblems-problemphyxmini0586-energyquantity}
  \lean{PhyXMiniProblems.ProblemPhyXMini0586.EnergyQuantity}
  A physical energy, carrying the dimension mass * length² / time²
\end{definition}
\begin{proof}
  This is the declared model definition of Energy Quantity.
\end{proof}
\begin{definition}[Wavelength Quantity]
  \label{def:physics:phyx-mini-0586:phyxminiproblems-problemphyxmini0586-wavelengthquantity}
  \lean{PhyXMiniProblems.ProblemPhyXMini0586.WavelengthQuantity}
  A physical wavelength, carrying the dimension of length
\end{definition}
\begin{proof}
  This is the declared model definition of Wavelength Quantity.
\end{proof}
\begin{definition}[energy In Electron Volts]
  \label{def:physics:phyx-mini-0586:phyxminiproblems-problemphyxmini0586-energyinelectronvolts}
  \lean{PhyXMiniProblems.ProblemPhyXMini0586.energyInElectronVolts}
  \uses{def:physics:phyx-mini-0586:phyxminiproblems-problemphyxmini0586-energyquantity}
  Read a physical energy as a real number of electron volts
\end{definition}
\begin{proof}
  This is the declared model definition of energy In Electron Volts.
\end{proof}
\begin{definition}[wavelength Readout]
  \label{def:physics:phyx-mini-0586:phyxminiproblems-problemphyxmini0586-wavelengthreadout}
  \lean{PhyXMiniProblems.ProblemPhyXMini0586.wavelengthReadout}
  \uses{def:physics:phyx-mini-0586:phyxminiproblems-problemphyxmini0586-wavelengthquantity}
  Read a physical wavelength as a real number in a selected length unit
\end{definition}
\begin{proof}
  This is the declared model definition of wavelength Readout.
\end{proof}
\begin{definition}[wavelength In Nanometers]
  \label{def:physics:phyx-mini-0586:phyxminiproblems-problemphyxmini0586-wavelengthinnanometers}
  \lean{PhyXMiniProblems.ProblemPhyXMini0586.wavelengthInNanometers}
  \uses{def:physics:phyx-mini-0586:phyxminiproblems-problemphyxmini0586-wavelengthquantity, def:physics:phyx-mini-0586:phyxminiproblems-problemphyxmini0586-wavelengthreadout}
  Nanometre readout of a physical wavelength
\end{definition}
\begin{proof}
  This is the declared model definition of wavelength In Nanometers.
\end{proof}
\begin{definition}[photon Energy Wavelength Constant Ev Nm]
  \label{def:physics:phyx-mini-0586:phyxminiproblems-problemphyxmini0586-photonenergywavelengthconstantevnm}
  \lean{PhyXMiniProblems.ProblemPhyXMini0586.photonEnergyWavelengthConstantEvNm}
  The rounded textbook value of h c used for photon-energy calculations when energy is measured in electron volts and wavelength in nanometres
\end{definition}
\begin{proof}
  This is the declared model definition of photon Energy Wavelength Constant Ev Nm.
\end{proof}
\begin{definition}[Particle Species]
  \label{def:physics:phyx-mini-0586:phyxminiproblems-problemphyxmini0586-particlespecies}
  \lean{PhyXMiniProblems.ProblemPhyXMini0586.ParticleSpecies}
  Particle species placed in the well
\end{definition}
\begin{proof}
  This is the declared model definition of Particle Species.
\end{proof}
\begin{definition}[Well Dimension]
  \label{def:physics:phyx-mini-0586:phyxminiproblems-problemphyxmini0586-welldimension}
  \lean{PhyXMiniProblems.ProblemPhyXMini0586.WellDimension}
  Number of spatial dimensions retained by the well model
\end{definition}
\begin{proof}
  This is the declared model definition of Well Dimension.
\end{proof}
\begin{definition}[Well Depth Kind]
  \label{def:physics:phyx-mini-0586:phyxminiproblems-problemphyxmini0586-welldepthkind}
  \lean{PhyXMiniProblems.ProblemPhyXMini0586.WellDepthKind}
  Whether the confining well has finite or idealized infinite depth
\end{definition}
\begin{proof}
  This is the declared model definition of Well Depth Kind.
\end{proof}
\begin{definition}[Energy Level Label]
  \label{def:physics:phyx-mini-0586:phyxminiproblems-problemphyxmini0586-energylevellabel}
  \lean{PhyXMiniProblems.ProblemPhyXMini0586.EnergyLevelLabel}
  The four energy labels printed beside the horizontal levels in the figure
\end{definition}
\begin{proof}
  This is the declared model definition of Energy Level Label.
\end{proof}
\begin{definition}[Absorption Feature]
  \label{def:physics:phyx-mini-0586:phyxminiproblems-problemphyxmini0586-absorptionfeature}
  \lean{PhyXMiniProblems.ProblemPhyXMini0586.AbsorptionFeature}
  The three wavelength features reported in the problem statement
\end{definition}
\begin{proof}
  This is the declared model definition of Absorption Feature.
\end{proof}
\begin{definition}[Spectral Region Kind]
  \label{def:physics:phyx-mini-0586:phyxminiproblems-problemphyxmini0586-spectralregionkind}
  \lean{PhyXMiniProblems.ProblemPhyXMini0586.SpectralRegionKind}
  The two kinds of spectral region distinguished by the figure
\end{definition}
\begin{proof}
  This is the declared model definition of Spectral Region Kind.
\end{proof}
\begin{definition}[Energy Level Figure]
  \label{def:physics:phyx-mini-0586:phyxminiproblems-problemphyxmini0586-energylevelfigure}
  \lean{PhyXMiniProblems.ProblemPhyXMini0586.EnergyLevelFigure}
  \uses{def:physics:phyx-mini-0586:phyxminiproblems-problemphyxmini0586-energylevellabel, def:physics:phyx-mini-0586:phyxminiproblems-problemphyxmini0586-spectralregionkind}
  Qualitative and schematic information transcribed from the primary raster. The vertical coordinates record only the displayed order; they are not energy measurements and are deliberately not assumed to be equally spaced
\end{definition}
\begin{proof}
  This is the declared model definition of Energy Level Figure.
\end{proof}
\begin{definition}[Finite Well Absorption Setup]
  \label{def:physics:phyx-mini-0586:phyxminiproblems-problemphyxmini0586-finitewellabsorptionsetup}
  \lean{PhyXMiniProblems.ProblemPhyXMini0586.FiniteWellAbsorptionSetup}
  \uses{def:physics:phyx-mini-0586:phyxminiproblems-problemphyxmini0586-energyquantity, def:physics:phyx-mini-0586:phyxminiproblems-problemphyxmini0586-wavelengthquantity, def:physics:phyx-mini-0586:phyxminiproblems-problemphyxmini0586-particlespecies, def:physics:phyx-mini-0586:phyxminiproblems-problemphyxmini0586-welldimension, def:physics:phyx-mini-0586:phyxminiproblems-problemphyxmini0586-welldepthkind, def:physics:phyx-mini-0586:phyxminiproblems-problemphyxmini0586-energylevellabel, def:physics:phyx-mini-0586:phyxminiproblems-problemphyxmini0586-absorptionfeature, def:physics:phyx-mini-0586:phyxminiproblems-problemphyxmini0586-energylevelfigure}
  Independent physical quantities and experimental predicates for the well. In particular, levelEnergy .E₂ is not assigned a numerical value here
\end{definition}
\begin{proof}
  This is the declared model definition of Finite Well Absorption Setup.
\end{proof}
\begin{definition}[Matches Finite One Dimensional Well Scenario]
  \label{def:physics:phyx-mini-0586:phyxminiproblems-problemphyxmini0586-matchesfiniteonedimensionalwellscenario}
  \lean{PhyXMiniProblems.ProblemPhyXMini0586.MatchesFiniteOneDimensionalWellScenario}
  \uses{def:physics:phyx-mini-0586:phyxminiproblems-problemphyxmini0586-finitewellabsorptionsetup}
  The qualitative finite-well scenario stated in the problem
\end{definition}
\begin{proof}
  This is the declared model definition of Matches Finite One Dimensional Well Scenario.
\end{proof}
\begin{definition}[Matches Problem Readouts]
  \label{def:physics:phyx-mini-0586:phyxminiproblems-problemphyxmini0586-matchesproblemreadouts}
  \lean{PhyXMiniProblems.ProblemPhyXMini0586.MatchesProblemReadouts}
  \uses{def:physics:phyx-mini-0586:phyxminiproblems-problemphyxmini0586-energyinelectronvolts, def:physics:phyx-mini-0586:phyxminiproblems-problemphyxmini0586-wavelengthinnanometers, def:physics:phyx-mini-0586:phyxminiproblems-problemphyxmini0586-finitewellabsorptionsetup}
  Numerical readouts stated in the prose: the continuum begins at 450.0 eV, the two line wavelengths are 14.588 nm and 4.8437 nm, and the continuum cutoff wavelength is 2.9108 nm. No energy for E₂ occurs in this premise
\end{definition}
\begin{proof}
  This is the declared model definition of Matches Problem Readouts.
\end{proof}
\begin{definition}[Matches Primary Energy Level Figure]
  \label{def:physics:phyx-mini-0586:phyxminiproblems-problemphyxmini0586-matchesprimaryenergylevelfigure}
  \lean{PhyXMiniProblems.ProblemPhyXMini0586.MatchesPrimaryEnergyLevelFigure}
  \uses{def:physics:phyx-mini-0586:phyxminiproblems-problemphyxmini0586-energylevelfigure}
  Primary-image transcription. The raster places E₁ < E₂ < E₃ < E₄, marks zero at the bottom of the energy axis, and begins the shaded nonquantized region at E₄. The auxiliary caption's claim of equal spacing is not used, because the primary image visibly has unequal schematic gaps
\end{definition}
\begin{proof}
  This is the declared model definition of Matches Primary Energy Level Figure.
\end{proof}
\begin{definition}[Has Physical Finite Well Parameters]
  \label{def:physics:phyx-mini-0586:phyxminiproblems-problemphyxmini0586-hasphysicalfinitewellparameters}
  \lean{PhyXMiniProblems.ProblemPhyXMini0586.HasPhysicalFiniteWellParameters}
  \uses{def:physics:phyx-mini-0586:phyxminiproblems-problemphyxmini0586-energyinelectronvolts, def:physics:phyx-mini-0586:phyxminiproblems-problemphyxmini0586-wavelengthinnanometers, def:physics:phyx-mini-0586:phyxminiproblems-problemphyxmini0586-finitewellabsorptionsetup}
  Positivity and ordering conditions selecting the physical finite-well regime
\end{definition}
\begin{proof}
  This is the declared model definition of Has Physical Finite Well Parameters.
\end{proof}
\begin{definition}[Satisfies Finite Well Photon Absorption Laws]
  \label{def:physics:phyx-mini-0586:phyxminiproblems-problemphyxmini0586-satisfiesfinitewellphotonabsorptionlaws}
  \lean{PhyXMiniProblems.ProblemPhyXMini0586.SatisfiesFiniteWellPhotonAbsorptionLaws}
  \uses{def:physics:phyx-mini-0586:phyxminiproblems-problemphyxmini0586-energyinelectronvolts, def:physics:phyx-mini-0586:phyxminiproblems-problemphyxmini0586-wavelengthinnanometers, def:physics:phyx-mini-0586:phyxminiproblems-problemphyxmini0586-photonenergywavelengthconstantevnm, def:physics:phyx-mini-0586:phyxminiproblems-problemphyxmini0586-finitewellabsorptionsetup}
  Photon absorption and ionization laws used by the solution: Eγ λ = h c ≈ 1240 eV·nm for every positive wavelength; absorption at λₐ raises the initially occupied level to E₂; absorption at λ\_b raises it to E₃; the cutoff photon reaches the continuum onset; every positive wavelength shorter than the cutoff is absorbable into the continuum. These relations contain no value for the first-excited-state energy and no answer choice
\end{definition}
\begin{proof}
  This is the declared model definition of Satisfies Finite Well Photon Absorption Laws.
\end{proof}
\begin{definition}[Answer Choice]
  \label{def:physics:phyx-mini-0586:phyxminiproblems-problemphyxmini0586-answerchoice}
  \lean{PhyXMiniProblems.ProblemPhyXMini0586.AnswerChoice}
  Labels of the four energy choices supplied with the problem
\end{definition}
\begin{proof}
  This is the declared model definition of Answer Choice.
\end{proof}
\begin{definition}[displayed Energy In Electron Volts]
  \label{def:physics:phyx-mini-0586:phyxminiproblems-problemphyxmini0586-displayedenergyinelectronvolts}
  \lean{PhyXMiniProblems.ProblemPhyXMini0586.displayedEnergyInElectronVolts}
  \uses{def:physics:phyx-mini-0586:phyxminiproblems-problemphyxmini0586-answerchoice}
  Electron-volt readout printed beside each answer label
\end{definition}
\begin{proof}
  This is the declared model definition of displayed Energy In Electron Volts.
\end{proof}
\begin{definition}[recorded Dataset Answer]
  \label{def:physics:phyx-mini-0586:phyxminiproblems-problemphyxmini0586-recordeddatasetanswer}
  \lean{PhyXMiniProblems.ProblemPhyXMini0586.recordedDatasetAnswer}
  \uses{def:physics:phyx-mini-0586:phyxminiproblems-problemphyxmini0586-answerchoice}
  The answer label recorded in the supplied dataset metadata
\end{definition}
\begin{proof}
  This is the declared model definition of recorded Dataset Answer.
\end{proof}
\begin{definition}[Rounds To Nearest Electron Volt]
  \label{def:physics:phyx-mini-0586:phyxminiproblems-problemphyxmini0586-roundstonearestelectronvolt}
  \lean{PhyXMiniProblems.ProblemPhyXMini0586.RoundsToNearestElectronVolt}
  A real energy readout rounds to a displayed whole number of electron volts
\end{definition}
\begin{proof}
  This is the declared model definition of Rounds To Nearest Electron Volt.
\end{proof}
\begin{definition}[Matches Displayed First Excited Energy]
  \label{def:physics:phyx-mini-0586:phyxminiproblems-problemphyxmini0586-matchesdisplayedfirstexcitedenergy}
  \lean{PhyXMiniProblems.ProblemPhyXMini0586.MatchesDisplayedFirstExcitedEnergy}
  \uses{def:physics:phyx-mini-0586:phyxminiproblems-problemphyxmini0586-energyinelectronvolts, def:physics:phyx-mini-0586:phyxminiproblems-problemphyxmini0586-finitewellabsorptionsetup, def:physics:phyx-mini-0586:phyxminiproblems-problemphyxmini0586-answerchoice, def:physics:phyx-mini-0586:phyxminiproblems-problemphyxmini0586-displayedenergyinelectronvolts, def:physics:phyx-mini-0586:phyxminiproblems-problemphyxmini0586-roundstonearestelectronvolt}
  A displayed choice agrees with the rounded first-excited-state energy
\end{definition}
\begin{proof}
  This is the declared model definition of Matches Displayed First Excited Energy.
\end{proof}
\begin{definition}[Is Unique Matching First Excited Energy]
  \label{def:physics:phyx-mini-0586:phyxminiproblems-problemphyxmini0586-isuniquematchingfirstexcitedenergy}
  \lean{PhyXMiniProblems.ProblemPhyXMini0586.IsUniqueMatchingFirstExcitedEnergy}
  \uses{def:physics:phyx-mini-0586:phyxminiproblems-problemphyxmini0586-finitewellabsorptionsetup, def:physics:phyx-mini-0586:phyxminiproblems-problemphyxmini0586-answerchoice, def:physics:phyx-mini-0586:phyxminiproblems-problemphyxmini0586-matchesdisplayedfirstexcitedenergy}
  A choice is the unique displayed value matching the modeled energy
\end{definition}
\begin{proof}
  This is the declared model definition of Is Unique Matching First Excited Energy.
\end{proof}
\begin{lemma}[ground State Energy Readout]
  \label{lem:physics:phyx-mini-0586:phyxminiproblems-problemphyxmini0586-groundstateenergyreadout}
  \lean{PhyXMiniProblems.ProblemPhyXMini0586.groundStateEnergyReadout}
  \uses{def:physics:phyx-mini-0586:phyxminiproblems-problemphyxmini0586-energyinelectronvolts, def:physics:phyx-mini-0586:phyxminiproblems-problemphyxmini0586-photonenergywavelengthconstantevnm, def:physics:phyx-mini-0586:phyxminiproblems-problemphyxmini0586-finitewellabsorptionsetup, def:physics:phyx-mini-0586:phyxminiproblems-problemphyxmini0586-matchesfiniteonedimensionalwellscenario, def:physics:phyx-mini-0586:phyxminiproblems-problemphyxmini0586-matchesproblemreadouts, def:physics:phyx-mini-0586:phyxminiproblems-problemphyxmini0586-hasphysicalfinitewellparameters, def:physics:phyx-mini-0586:phyxminiproblems-problemphyxmini0586-satisfiesfinitewellphotonabsorptionlaws}
  The continuum cutoff first determines the ground-state readout: the continuum begins at 450 eV, while a 2.9108 nm photon supplies the ionization gap
\end{lemma}
\begin{proof}
  The conclusion follows from the governing relations and figure readouts named in the statement of ground State Energy Readout.
\end{proof}
\begin{lemma}[first Excited Energy Readout]
  \label{lem:physics:phyx-mini-0586:phyxminiproblems-problemphyxmini0586-firstexcitedenergyreadout}
  \lean{PhyXMiniProblems.ProblemPhyXMini0586.firstExcitedEnergyReadout}
  \uses{def:physics:phyx-mini-0586:phyxminiproblems-problemphyxmini0586-energyinelectronvolts, def:physics:phyx-mini-0586:phyxminiproblems-problemphyxmini0586-photonenergywavelengthconstantevnm, def:physics:phyx-mini-0586:phyxminiproblems-problemphyxmini0586-finitewellabsorptionsetup, def:physics:phyx-mini-0586:phyxminiproblems-problemphyxmini0586-matchesfiniteonedimensionalwellscenario, def:physics:phyx-mini-0586:phyxminiproblems-problemphyxmini0586-matchesproblemreadouts, def:physics:phyx-mini-0586:phyxminiproblems-problemphyxmini0586-hasphysicalfinitewellparameters, def:physics:phyx-mini-0586:phyxminiproblems-problemphyxmini0586-satisfiesfinitewellphotonabsorptionlaws}
  Adding the 14.588 nm line-photon energy to the ground-state energy gives the unrounded first-excited-state energy readout
\end{lemma}
\begin{proof}
  The conclusion follows from the governing relations and figure readouts named in the statement of first Excited Energy Readout.
\end{proof}
% --- Archon declaration topology end ---
% --- Archon physics formalization source end ---
